\chapter{Parameterization}\label{ch:param}

A parameterization maps a surface onto a simple domain, usually a region of the plane. It is how
textures are painted on 3D models, and more generally it lets us process a surface with 2D
tools: many problems become easier in the plane. This chapter defines the problem, measures the
distortion of a map, and derives the classical linear methods (Tutte, harmonic and mean value
mappings) and a free-boundary conformal method (LSCM), following Tema 9 and the surveys of
\citet{Floater2005,Sheffer2006}.

\section{Applications}

\begin{itemize}
  \item \textbf{Texture mapping and 3D painting}: each vertex receives coordinates $(u,v)$ in a
  texture image, and the rendering pipeline interpolates them across the triangles; painting tools
  write directly into the texture through the inverse map.
  \item \textbf{Normal mapping}: simplification approximates the geometry with few triangles,
  but appearance depends on the normals. Decoupling geometry from appearance, the normals of the
  detailed mesh are stored in a texture of the simplified one, which then looks detailed.
  \item \textbf{Remeshing}: sampling a regular grid or a Delaunay triangulation in the parameter
  domain and mapping it back yields a new mesh (\cref{ch:remesh}).
  \item Surface fitting (splines over the parameter domain), editing, completion, interpolation,
  morphing and correspondence between shapes through a common domain, and visualisation.
\end{itemize}

\section{The problem}

Given a surface $S\subset\R^3$ and a domain $\Omega\subset\R^2$, find a bijective map
$\bu:S\to\Omega$, possibly subject to constraints such as fixed boundary positions
$\bu(\bx_i)=\bu_i$. For a triangle mesh, $\bu$ is piecewise linear, determined by the images
$\bu_i\in\R^2$ of the vertices, and bijectivity means that no triangle is flipped and that the
images do not overlap. Several difficulties appear immediately:
\begin{itemize}
  \item \textbf{Domain and topology}: a planar region can only receive surfaces with the topology of
  a disk. Closed surfaces or surfaces with handles must be cut along seams into one or several
  charts (an \emph{atlas}), or mapped to other domains such as the sphere or a coarse base mesh.
  \item \textbf{Boundary}: the boundary may be fixed to a convex polygon or left free.
  \item \textbf{Many solutions}: there are infinitely many bijections, so we must choose one that
  distorts the geometry as little as possible.
\end{itemize}

\section{Measuring distortion}

Restricted to a triangle, a piecewise linear map is affine: in an orthonormal frame of the triangle
plane its Jacobian $\mJ$ is a $2\times2$ matrix. Its singular values $\sigma_1\ge\sigma_2$ are the
maximal and minimal stretches: a small circle in the triangle maps to an ellipse with semi-axes
$\sigma_1$ and $\sigma_2$. A map is
\begin{itemize}
  \item \textbf{conformal} (angle preserving) if $\sigma_1=\sigma_2$;
  \item \textbf{equiareal} (area preserving) if $\sigma_1\sigma_2=1$;
  \item \textbf{isometric} (length preserving) if $\sigma_1=\sigma_2=1$, that is, both conformal and
  equiareal.
\end{itemize}
Isometric maps preserve intrinsic properties, those that depend only on angles and distances
measured on the surface, such as geodesic distances. Here the Theorema Egregium brings bad news:
Gaussian curvature is intrinsic, so a surface with $K\neq0$ has no isometric map to the plane. A
sphere cannot be flattened without distortion, which is why every map of the Earth compromises
(\cref{fig:projections}). The \emph{Mercator} projection
$(\lambda,\varphi)\mapsto(\lambda,\ln\tan(\tfrac\pi4+\tfrac\varphi2))$ is conformal (it maps
loxodromes, curves of constant bearing, to straight lines) but inflates areas towards the poles; the
\emph{Lambert} cylindrical projection $(\lambda,\varphi)\mapsto(\lambda,\sin\varphi)$ preserves
areas but shears angles. For meshes, the quantities $\sigma_1/\sigma_2\ge1$ (angle distortion) and
$\sigma_1\sigma_2$ (area distortion) per triangle measure the quality of a parameterization.

\fig[0.85\textwidth]{ch11_projections.pdf}{Two maps of the sphere with the same graticule and the
same circles of equal radius on the globe (Tissot indicatrices). Mercator keeps them circular but
enlarges them towards the poles; Lambert keeps their area but flattens them.}{fig:projections}

\section{The spring model and barycentric mappings}

Imagine a spring along every edge of the mesh, with stiffness $D_{ij}>0$. If the boundary vertices
are fixed in the plane, the interior vertices settle at the equilibrium that minimises the potential
energy
\begin{equation}
E=\frac12\sum_{(i,j)\in E}D_{ij}\norm{\bu_i-\bu_j}^2 .
\end{equation}
The fixed boundary prevents the degenerate solution in which all vertices collapse to a point.
Setting $\partial E/\partial\bu_i=\bzero$ for an interior vertex gives
\begin{equation}
\sum_{j\in\mathcal N(i)}D_{ij}(\bu_i-\bu_j)=\bzero\quad\Longleftrightarrow\quad
\bu_i=\sum_{j\in\mathcal N(i)}\lambda_{ij}\bu_j,\qquad \lambda_{ij}=\frac{D_{ij}}{\sum_kD_{ik}} :
\end{equation}
every interior vertex is a convex combination of its neighbours, with \emph{barycentric}
coordinates $\lambda_{ij}$. In practice:
\begin{enumerate}
  \item fix the boundary vertices on a convex polygon (triangle, square or circle), for example
  spaced proportionally to the boundary edge lengths;
  \item form the linear equations $\bu_i-\sum_j\lambda_{ij}\bu_j=0$ for the interior vertices,
  moving the known boundary terms to the right-hand side;
  \item solve two sparse linear systems with the same matrix, one for $u$ and one for $v$.
\end{enumerate}
In matrix form the system is $\mL\mathbf U=\mathbf B$. With the normalised weights $\lambda_{ij}$ the
matrix $\mL$ is not symmetric; using the raw weights $D_{ij}$, $L_{ij}=-D_{ij}$ and
$L_{ii}=\sum_jD_{ij}$, it is symmetric positive definite on the interior vertices and can be solved
with Cholesky or conjugate gradients.

\begin{theorem}[Maxwell--Tutte \citep{Tutte1963}, Floater \citep{Floater1997}]
If the mesh is a triangulated disk (a 3-connected planar graph), the boundary is mapped
bijectively onto a convex polygon and every interior vertex is a convex combination of its
neighbours with positive weights, then the resulting piecewise linear map is a valid embedding:
no triangle is flipped and the map is bijective.
\end{theorem}

Tutte's original \emph{barycentric drawing} uses uniform weights $D_{ij}=1$. It is always valid
but ignores the geometry: a planar mesh is not mapped to itself, so the method does not have
\emph{linear reproduction}. Weights proportional to the edge lengths do not fix this either.
Linear reproduction requires that each vertex's weights be barycentric coordinates of the vertex
with respect to its neighbours in the original (planar) position; for three neighbours they are
unique, but for more neighbours there are many possible choices, and the following two are the most
used.

\section{Conformal and harmonic maps}

The Riemann mapping theorem guarantees that any surface topologically equivalent to a disk can be
mapped conformally to the unit disk. In local isothermal coordinates a conformal map
$(u,v)$ satisfies the Cauchy--Riemann equations
\begin{equation}
\frac{\partial u}{\partial x}=\frac{\partial v}{\partial y},\qquad
\frac{\partial u}{\partial y}=-\frac{\partial v}{\partial x},
\end{equation}
which imply $\Lap u=\Lap v=0$: the coordinate functions of a conformal map are harmonic for the
Laplace--Beltrami operator. The converse does not hold, so isometric $\subset$ conformal $\subset$
harmonic. \emph{Harmonic maps} minimise the Dirichlet energy
$E_D(\bu)=\tfrac12\int_S\norm{\grad u}^2+\norm{\grad v}^2\dd A$ for given boundary values. They are
easier to compute than conformal maps (a linear problem) but do not necessarily preserve angles.

\begin{theorem}[Rad\'o--Kneser--Choquet]
If $\bu:S\to\R^2$ is harmonic and maps the boundary of $S$ homeomorphically onto the boundary of a
convex region $\Omega^\ast$, then $\bu$ is a bijection onto $\Omega^\ast$.
\end{theorem}

Discretising the Dirichlet energy with linear finite elements gives exactly the spring energy with
the cotangent weights of \cref{ch:smooth} \citep{Pinkall1993,Eck1995},
\begin{equation}
D_{ij}=\cot\alpha_{ij}+\cot\beta_{ij} .
\end{equation}
The \emph{discrete harmonic map} has linear reproduction and approximates the conformal map for
the given boundary. However, the cotangent weights are negative across non-Delaunay edges, so the
discrete theorem does not apply and the map is not guaranteed to be bijective: badly shaped meshes
can fold.

\paragraph{Mean value coordinates.} \citet{Floater2003} derived barycentric coordinates from the
mean value property of harmonic functions:
\begin{equation}
D_{ij}=\frac{\tan(\delta_{ij}/2)+\tan(\gamma_{ij}/2)}{\norm{\bx_i-\bx_j}},
\end{equation}
where $\delta_{ij}$ and $\gamma_{ij}$ are the angles at $\bx_i$ between the edge $ij$ and its two
neighbouring edges in the one-ring. The weights are always positive (the angles are smaller than
$\pi$), so the map is always valid by Tutte's theorem, they have linear reproduction, and the result
is visually very similar to the harmonic map, although not exactly conformal. Uniform and mean
value or harmonic weights with a fixed convex boundary are the most used methods in practice.

\begin{table}[tbp]
\centering\small
\begin{tabular}{@{}lccc@{}}
\toprule
Weights & Always bijective & Linear reproduction & Angle preservation \\
\midrule
Uniform (Tutte) & yes & no & poor \\
Cotangent (harmonic) & no (negative weights) & yes & good \\
Mean value & yes & yes & good \\
\bottomrule
\end{tabular}
\caption{Barycentric weights for fixed-boundary parameterization.}
\label{tab:weights}
\end{table}

\fig[0.85\textwidth]{ch11_tutte.png}{Fixed-boundary parameterizations of a face patch of the Max
Planck bust (3284 vertices, boundary on a circle) and the corresponding checkerboard textures.
Uniform weights ignore the geometry and give a median angle distortion $\sigma_1/\sigma_2=2.56$;
cotangent and mean value weights give $1.50$ and $1.52$. No triangle is flipped in any of the three
maps.}{fig:tutte}

\section{Free boundary: least squares conformal maps}

Fixing the boundary on a circle forces distortion near it: the boundary of the face patch is far
from circular, and the interior must absorb the difference. Free-boundary methods let the boundary
take its natural shape. \emph{Least squares conformal maps} \citep{Levy2002} write the
Cauchy--Riemann equations in complex form. In the local frame $(x,y)$ of a triangle $T$, with
$U=u+\mathrm iv$, a map is conformal if and only if
\begin{equation}
\frac{\partial U}{\partial x}+\mathrm i\,\frac{\partial U}{\partial y}=0 .
\end{equation}
On a triangle the derivatives are constant and linear in the unknown vertex coordinates, so the
\emph{conformal energy}
\begin{equation}
E_{\mathrm{LSCM}}(\mathbf U)=\sum_TA_T\left|\frac{\partial U}{\partial x}+\mathrm i\,\frac{\partial U}{\partial y}\right|^2
\end{equation}
is a quadratic form. It vanishes exactly for conformal maps, which exist only for developable
surfaces, so in general it is minimised in the least-squares sense. The energy is invariant under
similarity transformations of the plane (translation, rotation and scaling), so two vertices are
pinned to remove these degrees of freedom, and the minimiser is the solution of a sparse linear
least-squares problem. The same energy can be derived as the Dirichlet energy minus the area of the
image, the \emph{discrete conformal parameterization} of \citet{Desbrun2002}. With the boundary free,
the distortion drops dramatically (\cref{fig:lscm}).

\fig{ch11_free_boundary.png}{Fixed versus free boundary on the face patch. The harmonic map to a
circle has median $\log_2(\sigma_1/\sigma_2)=0.587$ (a ratio of $1.50$); LSCM, with the boundary
free, has $0.031$ (a ratio of $1.02$), and its checkerboard keeps square cells everywhere. Right:
histograms of the angle distortion per triangle.}{fig:lscm}

\paragraph{Other methods.} Angle-based flattening (ABF, ABF++ \citep{Sheffer2005}) optimises the
angles of the planar triangles directly and then reconstructs the positions. Local/global methods
\citep{Liu2008} fit a rotation (ARAP) or similarity (ASAP) per triangle and solve a global Poisson
system, alternating as in the deformation method of \cref{ch:deform}. For closed or high-genus
surfaces, the mesh is first cut into charts along seams chosen in high-curvature or hidden regions,
and the charts are packed into a texture atlas.

\begin{keypoints}
\begin{itemize}[leftmargin=*]
  \item A parameterization is a bijection from a disk-like surface to a planar domain; distortion is
  measured by the singular values of the per-triangle Jacobian.
  \item No isometric flattening exists when $K\neq0$ (Theorema Egregium): choose between angles
  (conformal) and areas (equiareal).
  \item Spring energies lead to barycentric mappings, $\bu_i=\sum_j\lambda_{ij}\bu_j$, solved as two
  sparse linear systems; with positive weights and a convex fixed boundary the map is bijective
  (Tutte, Floater).
  \item Cotangent weights discretise harmonic maps; mean value weights are always positive; LSCM
  frees the boundary and minimises the conformal energy with two pinned vertices.
\end{itemize}
\end{keypoints}

\section*{Exercises}
\begin{exercise}
Show that the singular values of the Jacobian of the Mercator projection at latitude $\varphi$ are
equal, and compute the area distortion as a function of $\varphi$.
\end{exercise}
\begin{exercise}
For a vertex with three neighbours in the plane, show that the barycentric coordinates are unique
and positive if the vertex lies inside the triangle of its neighbours.
\end{exercise}
\begin{exercise}
Prove that the conformal energy of LSCM is invariant under $U\mapsto aU+b$ with complex $a,b$.
Why are two pinned vertices enough to make the minimiser unique?
\end{exercise}
